\documentclass{article}
\usepackage{amsmath,hyperref}
\title{Shared admission for retained ordered warehouse classes}
\begin{document}
\maketitle
\paragraph{Interface and distinct plant.}
A previous loss leaves two mandatory classes \(7,6\) to repick. An old order at the native MRWS station must accept the classes \(7,6,5,4,3\) in that order. It has already accepted class \(8\). B holds \(5,4\); C holds \(3\). A second modeled battery fault before acceptance can disable either B or C. Select one replacement pattern for the old order at the dispatch cut, using a fresh named readiness report if present. In contrast to the radio suffix fixture, the preserved effect depends on existing physical class inventory and a native ordered station release, not a retransmitted keyed result or aging bandwidth. Staged stock needed for repicks is available by assumption. The shared cost budget counts new classes loaded during the horizon; each class costs one slot. A separate pick can stage an optional item. Time cuts, pickup slots, report delivery and assignment authority are authored; stock motion and sensing are not run.

\paragraph{Exact selection and transfer.}
Let \(b,c\in\{0,1\}\) request fresh replicas of the entire B and C batches respectively, and \(o\in\{0,1\}\) stage the optional item. The total budget use is \(2+2b+c+o\). In every history in the locally observed cell, a lost B requires \(b=1\), and a lost C requires \(c=1\). These conditions are necessary because the missing batch otherwise cannot reach the station through any legal mode in the catalogue. They suffice by serving the batches in order, staging at most two items per fresh batch on D and transferring before loading the next. Retained healthy goods not chosen are not claimed cleared. An unchanged admission interpreter enforces the same logical conditions with a pickup meter and requirement facts, checks exactly one old station completion and one separate side effect, and returns the intersection of feasible options over each cell. A direct finite enumeration ties. This is constructive reuse of the check, not an algorithmic advantage over the matched comparator.

\paragraph{Sharp capacity and report boundary.}
With one timely report, four slots suffice for every no-loss, B-loss and C-loss run; the B-loss branch proves four necessary. Under the same budget an optional item can be staged in the C-loss and no-loss branches but not the B-loss branch. When the report is absent at the cutoff, the observation cell contains both B-loss and C-loss. One pattern must then set both \(b=c=1\), requiring five slots, independently of whether it stages the optional item. A late authentic report cannot restore the preparation cut in this envelope. Serial full reconstruction is feasible at five slots but forgoes the smaller budget compared here. A model in which goods can be recovered from the failed carrier or the initial menu can adapt to a later read has a different conclusion.

\paragraph{Decoding and limitations.}
The pins reproduce the earlier MRWS source calibration: the native carrier inventory accepts each batch and the native station sends exactly one order completion. The selected fault callback identifies the disabled carrier; arrival at dispatch and truthful coverage of the native flag are fixture assumptions. Ordered consumption, actual staged inventory, and unrelated local unload run on the same fixture; the separate unload does not use the shared pickup budget. The native trace establishes selected callbacks and class-based acceptance, not a model of timed travel, unique provenance, or stock recovery under ROS or Unity. This is a second materially different development operation using the same admission interpreter, not a held-out test, integrated I1 discharge, or a claim of superiority to a source-aware planner.
\end{document}
